% ============================================================================
% DROP-IN REPLACEMENT FOR YOUR CHESS DIAGRAM
% ============================================================================
% This replaces the chess minipage in your original figure.
% It adds an explicit dependency graph showing the BFS + minimax structure.
%
% Required packages (add to your preamble if not already there):
% \usepackage{tikz}
% \usetikzlibrary{positioning,arrows.meta}
% \usepackage{tcolorbox}
% \usepackage{skak}  % for chess symbols
% \usepackage{xcolor}
% ============================================================================

    \begin{minipage}{0.48\textwidth}
        \begin{tcolorbox}[
            title=\textbf{Chess: Minimax Search} \hfill {\small\colorbox{orange!20}{Easy}},
            colframe=orange!70!black,
            colback=orange!3,
            boxrule=0.8pt
        ]
        \footnotesize

        \textbf{Dependency graph:}
        \begin{center}
        \begin{tikzpicture}[
            node distance=0.35cm and 0.4cm,
            every node/.style={font=\tiny},
            bfs/.style={draw, rounded corners=1.5pt, fill=orange!12, minimum width=0.55cm, minimum height=0.35cm, thick, draw=orange!45},
            ord/.style={draw, rounded corners=1.5pt, fill=orange!28, minimum width=0.85cm, minimum height=0.35cm, thick, draw=orange!55},
            outputnode/.style={draw, rounded corners=1.5pt, fill=orange!48, minimum width=0.7cm, minimum height=0.35cm, thick, draw=orange!70, font=\tiny\bfseries},
            arrow/.style={-{Stealth[length=1.2mm]}, semithick, orange!55!black}
        ]
            % BFS layer
            \node[bfs] (b1) {$d_{K1}$};
            \node[bfs, right=0.15cm of b1] (b2) {$d_{K2}$};
            \node[bfs, right=0.15cm of b2] (b3) {$d_{K3}$};
            \node[bfs, right=0.15cm of b3] (b4) {$d_{12}$};
            \node[bfs, right=0.15cm of b4] (b5) {$d_{13}$};
            \node[bfs, right=0.15cm of b5] (b6) {$d_{23}$};

            % Orderings layer (3! = 6)
            \node[ord, below=0.5cm of b1, xshift=0.3cm] (o1) {1-2-3};
            \node[ord, right=0.1cm of o1] (o2) {1-3-2};
            \node[ord, right=0.1cm of o2] (o3) {2-1-3};
            \node[ord, below=0.25cm of o1] (o4) {2-3-1};
            \node[ord, right=0.1cm of o4] (o5) {3-1-2};
            \node[ord, right=0.1cm of o5] (o6) {3-2-1};

            % Output
            \node[outputnode, below=0.45cm of o5] (out) {minimax};

            % Arrows BFS to orderings (simplified)
            \draw[arrow] (b2.south) -- ++(0,-0.12) -| (o1.north);
            \draw[arrow] (b4.south) -- ++(0,-0.08) -| (o2.north);
            \draw[arrow] (b6.south) -- ++(0,-0.12) -| (o3.north);

            % Arrows orderings to output
            \draw[arrow] (o4.south) -- (out.north west);
            \draw[arrow] (o5.south) -- (out.north);
            \draw[arrow] (o6.south) -- (out.north east);
        \end{tikzpicture}
        \end{center}

        \vspace{-0.15cm}
        {\scriptsize BFS computes knight distances; minimax evaluates $3!$ capture orderings (Alice max, Bob min).}

        \vspace{0.15cm}
        \hfill \textbf{Answer:} [Redacted]
        \end{tcolorbox}
    \end{minipage}%

% ============================================================================
% WHAT THIS SHOWS:
% - Top row: 6 BFS computations (distances between knight K and pawns 1,2,3)
% - Middle rows: All 3! = 6 possible capture orderings
% - Bottom: Minimax aggregation (Alice maximizes, Bob minimizes)
%
% The dependency structure is:
%   BFS distances → evaluate each ordering → minimax over all orderings
% ============================================================================
